\section{两个线性映射的张量积}

记号约定:

\begin{itemize}
	\item $A$是环,$E_{1},E_{2},E_{3}$是右$A$-模,$F_{1},F_{2},F_{3}$是左$A$-模.
	\item $f\in\Hom_{A}\left(E_{1},E_{2}\right),f^{\prime}\in\Hom_{A}\left(E_{2},E_{3}\right),g\in\Hom_{A}\left(F_{1},F_{2}\right),g^{\prime}\in\Hom_{A}\left(F_{2},F_{3}\right)$.
\end{itemize}

本小节中,$E_{1}$,$E_{2}$是右$A$-模,$F_{1}$,$F_{2}$是左$A$-模,$f:E_{1}\to E_{2}$,$g:F_{1}\to F_{2}$是模同态.

\begin{definition}{线性映射的张量积}{}
	我们称$f\otimes_{\Z}g:E_{1}\otimes_{A}F_{1}\to E_{2}\otimes_{A}F_{2},x\otimes y\mapsto f\left(x\right)\otimes g\left(y\right)$为$f$和$g$的张量积.
\end{definition}

\begin{proposition}{}{}
	\begin{enumerate}
		\item $\Hom_{A}\left(E_{1},E_{2}\right)\times\Hom_{A}\left(F_{1},F_{2}\right)\to\Hom_{\Z}\left(E_{1}\otimes_{A}F_{1},E_{2}\otimes_{A}F_{2}\right),\left(u,v\right)\mapsto u\otimes_{Z}v$是$\Z$-双线性映射.
		\item (1)中的映射诱导如下的典范$\Z$-线性映射
		\begin{gather*}
			\Hom_{A}\left(E_{1},E_{2}\right)\otimes_{\Z}\Hom_{A}\left(F_{1},F_{2}\right)\to\Hom_{\Z}\left(E_{1}\otimes_{A}F_{1},E_{2}\otimes_{A}F_{2}\right),\\
			u\otimes_{\Z}v\mapsto u\otimes_{A}v.
		\end{gather*}
	\end{enumerate}
\end{proposition}

\begin{proposition}{张量积的函子性}{}
	如下图表交换.
	\begin{center}
		\begin{tikzcd}
			E_{1}\otimes_{A}F_{1} \arrow[dd, "\left(f^{\prime}\circ f\right)\otimes\left(g^{\prime}\circ g\right)"'] \arrow[rr, "f\otimes g"] &  & E_{2}\otimes_{A}F_{2} \arrow[dd, "f^{\prime}\otimes g^{\prime}"] \\
			&  &                                                                \\
			E_{3}\otimes_{A}F_{3} \arrow[rr, "\id_{E_{3}\otimes_{A}F_{3}}", equals]                                                                           &  & E_{3}\otimes_{A}F_{3}                                         
		\end{tikzcd}
	\end{center}
\end{proposition}
